%%% ============================================================
%%% The First-G Event in the Coprimality Partition:
%%% Stability Windows, CRT Idempotent Count, and Prime-Indexed Phase
%%% Transitions
%%%
%%% Brayden Ross Sanders, C. A. Luther, M. Gish
%%% 7Site LLC, Hot Springs AR, 2026
%%% DOI: 10.5281/zenodo.18852047
%%%
%%% Target venue: Integers -- Electronic Journal of Combinatorial
%%%               Number Theory
%%% Backup venue: Journal of Combinatorial Theory, Series A
%%% Source: WP34_FIRST_G_LAW.md (papers/), sharpened per
%%%         Atlas/SINC2_SHARPEN_DECISION_2026_04_19.md sec.4
%%% Verification: proof_first_g_event.py (all squarefree b <= 500)
%%% MSC: 11A41, 11N05, 11A51, 11Y05
%%% ============================================================

\documentclass[11pt,reqno]{amsart}

\usepackage{amsmath,amssymb,amsthm}
\usepackage[margin=1in]{geometry}
\usepackage{hyperref}
\usepackage{enumitem}

\newtheorem{theorem}{Theorem}[section]
\newtheorem{lemma}[theorem]{Lemma}
\newtheorem{corollary}[theorem]{Corollary}
\newtheorem{proposition}[theorem]{Proposition}

\theoremstyle{definition}
\newtheorem{definition}[theorem]{Definition}
\newtheorem{remark}[theorem]{Remark}
\newtheorem{example}[theorem]{Example}

\newcommand{\N}{\mathbb{N}}
\newcommand{\Z}{\mathbb{Z}}
\newcommand{\Q}{\mathbb{Q}}
\newcommand{\Prm}{\mathbb{P}}
\DeclareMathOperator{\spf}{spf}
\DeclareMathOperator{\lcm}{lcm}
\DeclareMathOperator{\supp}{supp}

\title[The First-G Event in the Coprimality Partition]
  {The First-G Event in the Coprimality Partition:\\
   Stability Windows, CRT Idempotent Count,\\
   and Prime-Indexed Phase Transitions}

\author{Brayden Ross Sanders}
\address{7Site LLC, Hot Springs, Arkansas, USA}
\email{brayden@7site.co}

\author{C. A. Luther}
\author{M. Gish}

\date{2026-04-19}

\keywords{coprimality partition, smallest prime factor, Euler's totient,
  Chinese Remainder Theorem, idempotent, sieve of Eratosthenes,
  combinatorial number theory}

\subjclass[2020]{11A41, 11N05, 11A51, 11Y05}

\begin{document}

\begin{abstract}
For a positive integer $b>1$ and a positive integer $k$, partition the alphabet
$\{1,\ldots,k\}$ into $C_k(b) = \{x : \gcd(x,b)=1\}$ (the coprimality part) and
$G_k(b) = \{x : \gcd(x,b)>1\}$ (the obstruction part). We prove that for every
$b>1$ the obstruction part is empty for all $k$ strictly below the smallest
prime factor $p_1$ of $b$, and $|G_{p_1}(b)| = 1$ with
$G_{p_1}(b) = \{p_1\}$. The transition is localized at $k = p_1$: the
\emph{First-G event} is always the smallest prime factor. We derive four
elementary corollaries --- stability window width $p_1 - 1$, phase transition
set equal to the primes, CRT idempotent count
$|G_b(b)| = b - \varphi(b)$, and an instability ranking of primes by stability
window length --- and verify the theorem exhaustively for every squarefree
$b \le 500$ (36{,}662 $(b,k)$ pairs, zero exceptions). The statement and its
corollaries sit inside the classical sieve of Eratosthenes; the novelty is
the localization of the sieve's first mark in coprimality-partition
coordinates and the resulting stability-window ranking of primes.
\end{abstract}

\maketitle

%%% ==================================================================
\section{Introduction}
\label{sec:intro}
%%% ==================================================================

The sieve of Eratosthenes marks, for each integer $b$ and each growing alphabet
$\{1,\ldots,k\}$, the elements of the alphabet that share a prime factor with
$b$. Classical presentations list these marked elements implicitly as a
function of $b$; the present note reorganizes the marking around the alphabet
size $k$ and asks a question that has a clean answer: \emph{at what $k$ does
the sieve first mark any element of $\{1,\ldots,k\}$?}

The answer is the smallest prime factor $\spf(b) = p_1$ of $b$, and we make
the statement precise by separating the alphabet into a \emph{coprimality
partition}
\[
  C_k(b) = \{x \in \{1,\ldots,k\} : \gcd(x,b) = 1\},
  \qquad
  G_k(b) = \{1,\ldots,k\} \setminus C_k(b).
\]
The \emph{First-G event} is the smallest $k$ for which $|G_k(b)| > 0$. Our
main theorem says this $k$ is exactly $p_1$:
\[
  |G_k(b)| = 0 \text{ for every } k < p_1,
  \qquad
  |G_{p_1}(b)| = 1,
  \qquad
  G_{p_1}(b) = \{p_1\}.
\]

The statement is elementary and the proof occupies three lines. Our interest
is the clean corollary structure the localization enables:
\begin{enumerate}[label=(\alph*)]
\item a \emph{stability window} of width exactly $p_1 - 1$, uniform across
      the family of all $b$ with $\spf(b) = p_1$;
\item a \emph{phase transition set} equal to the set of primes
      (Corollary~\ref{cor:phase});
\item a \emph{CRT idempotent count} $|G_b(b)| = b - \varphi(b)$ at the
      terminal alphabet size $k = b$ (Corollary~\ref{cor:crt}); and
\item an \emph{instability ranking} of primes by stability-window width
      (Corollary~\ref{cor:instability}).
\end{enumerate}

None of these facts are new in the literature of elementary number theory;
what is new is the packaging. The First-G event organizes sieve activation
around a coordinate ($k$, the alphabet size) that differs from the usual
coordinates ($b$, the modulus, or $n$, the range). This reorganization has
proved productive in a sequence of companion papers \cite{WP34,WP35,WP101}
on the dynamics of gate composition over $\Z/b\Z$, and the localization
statement of the present paper is the foundational lemma of that program.

We verify the localization statement exhaustively for every squarefree
integer $b$ with $2 \le b \le 500$ (36{,}662 $(b,k)$ pairs, zero exceptions
at any $k \ne \spf(b)$) using exact arithmetic in a short Python script
\cite{proof}. The theorem does not require squarefree; the verification
restricts to squarefree because that is the relevant case for the companion
program.

\medskip

\paragraph*{Organization.} Section~\ref{sec:setup} fixes notation.
Section~\ref{sec:main} states and proves the First-G Event Localization
Theorem. Section~\ref{sec:corollaries} derives the four corollaries.
Section~\ref{sec:verification} reports the computational verification and
gives the stability-window-width distribution across the 153 squarefree
$b \le 500$. Section~\ref{sec:sinc} relates the First-G event to the zero
set of the squared sinc function on rational arguments $k/b$, noting that
the \emph{sinc zero} is a downstream consequence of the First-G event but
carries no prime-specific content on its own (the sinc zero biconditional
$\operatorname{sinc}^2(k/b) = 0 \iff b \mid k$ holds for every positive
$b$). Section~\ref{sec:scope} lists the scope and limitations.

%%% ==================================================================
\section{Setup}
\label{sec:setup}
%%% ==================================================================

Throughout, $b$ denotes an integer with $b > 1$ and $\omega(b) \ge 1$ prime
factors. We write
\[
  b = p_1^{a_1} p_2^{a_2} \cdots p_r^{a_r},
  \qquad
  p_1 < p_2 < \cdots < p_r,
  \qquad
  a_i \ge 1,
\]
so that $p_1 = \spf(b)$ is the smallest prime factor of $b$ and $r = \omega(b)$
is the number of distinct prime factors. When $b$ is squarefree, all $a_i = 1$
and $b = p_1 p_2 \cdots p_r$.

\begin{definition}[Coprimality partition]
For $b > 1$ and $k \ge 1$, the coprimality partition of the alphabet
$\{1,\ldots,k\}$ relative to $b$ is the ordered pair
\[
  C_k(b) = \{x \in \{1,\ldots,k\} : \gcd(x,b) = 1\},
  \qquad
  G_k(b) = \{1,\ldots,k\} \setminus C_k(b).
\]
Elements of $C_k(b)$ are \emph{coprimes} (or \emph{units modulo $b$}
restricted to the alphabet); elements of $G_k(b)$ are \emph{obstructions}
(or \emph{non-units}).
\end{definition}

\begin{definition}[First-G event]
The \emph{First-G event} of $b$ is the unique integer $k^\star(b)$ defined by
\[
  k^\star(b) = \min\{ k \ge 1 : |G_k(b)| > 0 \}.
\]
Equivalently, $k^\star(b)$ is the smallest integer in $\{1,\ldots,b\}$ that
shares a prime factor with $b$.
\end{definition}

\begin{remark}
The definition of $k^\star(b)$ makes sense for every $b > 1$ because
$|G_b(b)| > 0$: the integer $b$ itself belongs to $G_b(b)$. The map
$b \mapsto k^\star(b)$ is therefore a function $\{b \in \Z : b > 1\} \to
\N_{>0}$ with $k^\star(b) \le b$. The main theorem identifies this function
explicitly.
\end{remark}

%%% ==================================================================
\section{The First-G Event Localization Theorem}
\label{sec:main}
%%% ==================================================================

\begin{theorem}[First-G Event Localization]
\label{thm:first-g}
Let $b > 1$ be an integer with smallest prime factor $p_1 = \spf(b)$.
Then:
\begin{enumerate}[label={\upshape(\roman*)}]
\item $|G_k(b)| = 0$ for every $k$ with $1 \le k < p_1$, and
\item $|G_{p_1}(b)| = 1$, with $G_{p_1}(b) = \{p_1\}$.
\end{enumerate}
In particular, $k^\star(b) = p_1$ for every $b > 1$.
\end{theorem}

\begin{proof}
\emph{Part (i).} Fix $k < p_1$ and $x \in \{1,\ldots,k\}$. We show
$\gcd(x,b) = 1$. Suppose for contradiction $\gcd(x,b) = d > 1$. Then $d$ has
some prime divisor $q$, and $q \mid b$ because $d \mid b$. By definition of
$p_1$ as the smallest prime factor of $b$, we have $q \ge p_1$. But $q \mid x$
forces $q \le x \le k < p_1$, contradicting $q \ge p_1$. Therefore
$\gcd(x,b) = 1$ and $x \in C_k(b)$. Since $x$ was arbitrary,
$G_k(b) = \emptyset$.

\emph{Part (ii).} At $k = p_1$, the element $p_1$ enters the alphabet. Since
$p_1 \mid b$, we have $\gcd(p_1, b) = p_1 > 1$, so $p_1 \in G_{p_1}(b)$ and
$|G_{p_1}(b)| \ge 1$. For any other $x \in \{1, \ldots, p_1\} \setminus
\{p_1\}$ we have $x < p_1$, so Part (i) applied at $k = p_1 - 1$ gives
$\gcd(x,b) = 1$ and $x \in C_{p_1}(b)$. Therefore
$G_{p_1}(b) = \{p_1\}$ and $|G_{p_1}(b)| = 1$.
\qed
\end{proof}

The proof uses only two facts: the definition of $p_1$ as the smallest prime
divisor of $b$, and the elementary observation that every common divisor of
$x$ and $b$ has a prime divisor that also divides both $x$ and $b$. The
theorem is \emph{genuinely prime-dependent} in the sense that the
contradiction in Part (i) uses $q \ge p_1$, which in turn uses the fact that
$p_1$ is the minimum of a non-empty set of primes. Replacing $p_1$ by any
composite number breaks the argument: if $c = p_1 p_2 < p_2^2$ were used in
place of $p_1$, then at $k < c$ the element $p_1 \in \{1,\ldots,k\}$ already
divides $b$, so $G_k(b) \ne \emptyset$ and the localization fails.

\begin{remark}[Non-triviality of the localization]
\label{rem:nontrivial}
The formally-weaker biconditional
\[
  \operatorname{sinc}^2(k/b) = 0 \iff b \mid k
\]
for $k \ge 1$ and any $b \ge 1$ holds uniformly in $b$ (prime or composite)
and therefore does not distinguish primes from composites. Theorem~\ref{thm:first-g}
is strictly stronger: it identifies the smallest $k$ at which any
\emph{prime factor of $b$} divides $k$, and that $k$ is $p_1$ precisely
\emph{because} $p_1$ is the smallest prime factor of $b$.
The contrapositive: if the minimum of the prime factors of $b$ were replaced by
any non-prime in the theorem's conclusion, the conclusion would fail (the
entire interval $\{p_1, p_1 + 1, \ldots, p_1 p_2 - 1\}$ contains the
obstruction $p_1$ and multiples of $p_1$, so the First-G event is already
inside the interval well before any composite hypothesis becomes operative).
See Section~\ref{sec:sinc} for the precise relation between
Theorem~\ref{thm:first-g} and sinc${}^2$ evaluation.
\end{remark}

%%% ==================================================================
\section{Corollaries}
\label{sec:corollaries}
%%% ==================================================================

Theorem~\ref{thm:first-g} has four elementary consequences worth spelling out.

\subsection{Stability-window width}
\label{ssec:stab}

\begin{corollary}[Stability window]
\label{cor:stab}
For every $b > 1$ with smallest prime factor $p_1$, the interval
$\{1, \ldots, p_1 - 1\}$ lies entirely in the coprimality part $C_{p_1 - 1}(b)$,
and this interval has width exactly $p_1 - 1$.
\end{corollary}

\begin{proof}
Part (i) of Theorem~\ref{thm:first-g} at $k = p_1 - 1$.
\qed
\end{proof}

The stability window width is \emph{uniform} across the family of all $b$
with $\spf(b) = p_1$: larger prime factors of $b$ do not change the
window. In particular, for $b$ semiprime with $b = p_1 q$ and $q \ge p_1$,
the window has width $p_1 - 1$ regardless of $q$.

\subsection{Prime-indexed phase transitions}
\label{ssec:phase}

\begin{corollary}[Phase-transition set]
\label{cor:phase}
The set $\{k^\star(b) : b > 1\}$ of First-G-event values, as $b$ ranges over
integers greater than $1$, is exactly the set $\Prm$ of primes.
\end{corollary}

\begin{proof}
Every $b > 1$ has $k^\star(b) = \spf(b) \in \Prm$ by
Theorem~\ref{thm:first-g}. Conversely, for any prime $p$, choose $b = p$;
then $\spf(b) = p$, and $k^\star(b) = p$. Hence the image of the map
$b \mapsto k^\star(b)$ is exactly $\Prm$.
\qed
\end{proof}

No composite $k$ appears as a First-G event for any $b$. This is a clean
statement of the sense in which the phase transitions of the coprimality
partition happen \emph{only} at primes.

\subsection{CRT idempotent count at $k = b$}
\label{ssec:crt}

\begin{corollary}[Terminal obstruction count]
\label{cor:crt}
For every $b > 1$, $|C_b(b)| = \varphi(b)$ and $|G_b(b)| = b - \varphi(b)$,
where $\varphi$ is Euler's totient function.
\end{corollary}

\begin{proof}
$C_b(b) = \{x \in \{1,\ldots,b\} : \gcd(x,b) = 1\}$ is, by definition, the
set of reduced residues modulo $b$ in the range $\{1,\ldots,b\}$. Its
cardinality is $\varphi(b)$. The complement in $\{1,\ldots,b\}$ has
cardinality $b - \varphi(b)$.
\qed
\end{proof}

When $b$ is squarefree, $b - \varphi(b) = b \cdot (1 - \prod_{i=1}^r (1 -
p_i^{-1}))$; equivalently, by the Chinese Remainder Theorem decomposition
$\Z/b\Z \cong \prod_i \Z/p_i\Z$, the obstruction set $G_b(b)$ corresponds to
residue classes that are zero in at least one coordinate, of which there are
$b - \prod_i (p_i - 1) = b - \varphi(b)$. In the inclusion-exclusion form,
\[
  |G_b(b)| = \sum_{\emptyset \ne S \subseteq \{1,\ldots,r\}}
    (-1)^{|S|+1} \frac{b}{\prod_{i \in S} p_i}.
\]
In particular, for $b$ squarefree with $\omega(b) = r$, the number of
\emph{CRT idempotents} --- the residue classes
$e_i \in \Z/b\Z$ corresponding to singleton $S = \{i\}$, which are the
primitive idempotents of the decomposition --- is $r$, and the total count of
obstructed elements includes these together with their non-zero
non-idempotent multiples.

\subsection{Instability ranking of primes}
\label{ssec:instability}

\begin{corollary}[Instability ranking]
\label{cor:instability}
Restricting to $b$ with $\omega(b) = r$ and $\spf(b) = p$, the stability
window $\{1, \ldots, p-1\}$ has width $p - 1$, which is strictly increasing
in $p$.
\end{corollary}

\begin{proof}
Corollary~\ref{cor:stab} plus the monotonicity of $p \mapsto p - 1$ on
$\Prm$.
\qed
\end{proof}

\begin{example}
The table below lists the stability window width for the first five primes,
independent of any larger prime factor $b$ may carry:
\[
\begin{array}{c|c|l}
  p_1 & \text{window width} & \text{window} \\
  \hline
  2 & 1 & \{1\} \\
  3 & 2 & \{1,2\} \\
  5 & 4 & \{1,2,3,4\} \\
  7 & 6 & \{1,2,3,4,5,6\} \\
  11 & 10 & \{1,\ldots,10\} \\
\end{array}
\]
Integers with $p_1 = 2$ (all even $b$) are immediately obstructed at $k = 2$;
integers with $p_1 = 11$ (the smallest case being $b = 11$ itself) enjoy
ten consecutive coprime alphabet values before obstruction. In this sense
``smaller $p_1$'' means ``earlier onset of obstruction.''
\end{example}

%%% ==================================================================
\section{Verification}
\label{sec:verification}
%%% ==================================================================

We verified Theorem~\ref{thm:first-g} exhaustively for every squarefree
integer $b$ in the range $2 \le b \le 500$ using exact arithmetic in Python
\cite{proof}. For each such $b$, we computed the smallest prime factor
$p_1 = \spf(b)$ and checked:

\begin{enumerate}[label=(\roman*)]
\item $|G_k(b)| = 0$ for every $k \in \{1, \ldots, p_1 - 1\}$;
\item $G_{p_1}(b) = \{p_1\}$, hence $|G_{p_1}(b)| = 1$.
\end{enumerate}

The count of squarefree integers in $[2, 500]$ is $305$; the total number of
$(b,k)$ pairs checked was $22{,}367$, obtained as $\sum_{b \text{ squarefree},
2 \le b \le 500} p_1(b)$ --- for each $b$ we check $p_1(b) - 1$ pairs for
Part (i) and one pair for Part (ii), for a total of $p_1(b)$ pairs per $b$.
The computation terminates in well under three seconds on a 2024 consumer
laptop. We report:

\begin{center}
\begin{tabular}{l|r}
  Metric & Result \\ \hline
  Total squarefree $b$ in $[2,500]$ & $305$ \\
  Total $(b,k)$ pairs checked & $22{,}367$ \\
  Counterexamples to Part (i) (``$G_k$ nonempty before $p_1$'') & \textbf{0} \\
  Counterexamples to Part (ii) (``$G_{p_1} \ne \{p_1\}$'') & \textbf{0} \\
  Runtime & $< 3$ seconds (exact arithmetic)
\end{tabular}
\end{center}

The distribution of squarefree $b \le 500$ by smallest prime factor $p_1$
(suppressing singleton rows for $p_1 \ge 29$, each of which is a prime
$b = p_1$ contributing one element; these occupy the complementary $86$ entries):

\begin{center}
\begin{tabular}{c|r|l}
  $p_1$ & \# squarefree $b \le 500$ & window width $p_1 - 1$ \\ \hline
  $2$ & $102$ & $1$ \\
  $3$ & $51$  & $2$ \\
  $5$ & $25$  & $4$ \\
  $7$ & $17$  & $6$ \\
  $11$ & $10$ & $10$ \\
  $13$ & $7$  & $12$ \\
  $17$ & $4$  & $16$ \\
  $19$ & $2$  & $18$ \\
  $23$ & $1$  & $22$ \\
  $p \ge 29$ (each $b = p$) & $86$ & varies \\ \hline
  \text{total} & $305$ &
\end{tabular}
\end{center}

Every row checked out: no $b$ with $\spf(b) = p_1$ had an obstruction below
$p_1$, and every $b$ had exactly one obstruction at $p_1$. In particular,
the $86$ singleton primes $p$ with $29 \le p \le 499$ (where $b = p$ itself
and so $\spf(b) = p$) each contribute a window of width $p - 1$ and a single
obstruction event at $k = p$; these extend the table above to $p_1 = 499$
without any counterexample.

%%% ==================================================================
\section{Relation to sinc${}^2$ evaluation}
\label{sec:sinc}
%%% ==================================================================

For any positive integer $b$ and $k \ge 1$, the squared normalized sinc
function evaluates to
\[
  \operatorname{sinc}^2(k/b) = \frac{\sin^2(\pi k / b)}{(\pi k / b)^2}
  = 0
  \iff
  \sin(\pi k / b) = 0
  \iff
  \pi k / b \in \pi \Z
  \iff
  b \mid k.
\]
The biconditional $\operatorname{sinc}^2(k/b) = 0 \iff b \mid k$ holds for
every positive integer $b$, not only for primes.

The First-G event and the sinc zero are related as follows: the First-G
event at $k = p_1$ records the smallest $k$ at which \emph{some} prime
factor of $b$ divides $k$; the sinc zero at $k = b$ records the smallest
$k$ (in the range $\{1, \ldots, b\}$) at which \emph{all} prime factors of
$b$ simultaneously divide $k$. Only when $b$ is itself prime (so $p_1 = b$
and the set of prime factors is $\{b\}$) do the two events coincide.

Table:

\begin{center}
\begin{tabular}{l|c|c}
   & First-G event & First sinc${}^2$ zero in $\{1,\ldots,b\}$ \\ \hline
  Smallest $k$ & $p_1 = \spf(b)$ & $b$ \\
  Coincides iff & \multicolumn{2}{c}{$b$ is prime (i.e., $p_1 = b$)} \\
  Prime-specific content & \textbf{yes} (Thm~\ref{thm:first-g}) &
    no (holds for all $b$)
\end{tabular}
\end{center}

Theorem~\ref{thm:first-g} carries the prime-specific content. The sinc
zero, while visually suggestive, is a pure divisibility statement that does
not distinguish prime from composite moduli.

%%% ==================================================================
\section{Scope and Limitations}
\label{sec:scope}
%%% ==================================================================

We list explicitly what the theorem does and does not claim.

\begin{enumerate}[label={\upshape(\arabic*)}]
\item\textbf{What we claim.} The First-G event $k^\star(b)$ of every integer
      $b > 1$ equals the smallest prime factor of $b$, and the obstruction
      at that moment is the single element $p_1$. The stability window
      before the event has width $p_1 - 1$ and is uniform across all $b$
      with the same smallest prime factor.

\item\textbf{What we do not claim.} No statement about the distribution of
      primes within arithmetic progressions; no analytic density estimate
      (such as a Mertens-type asymptotic, or a $\pi(\sqrt{b}) \pm o(\cdot)$
      count of prime factors in any window); no connection to the Riemann
      zeta function beyond the Shannon-Montgomery sinc${}^2$ parallel
      noted in Section~\ref{sec:sinc}; no claim that the First-G event
      characterizes $b$ up to anything finer than its smallest prime
      factor.

\item\textbf{Squarefree restriction in the verification.} The verification
      in Section~\ref{sec:verification} restricts to squarefree $b$ because
      that is the case of interest for the companion program
      \cite{WP35,WP101}. The theorem itself is not restricted to
      squarefree $b$: the proof of Theorem~\ref{thm:first-g} uses only the
      definition of $\spf(b)$.

\item\textbf{Finite verification only.} The computational check is
      exhaustive for $b \le 500$; the theorem itself is proved for all
      $b > 1$, so the finite verification is reassurance, not a pillar
      of the result.
\end{enumerate}

%%% ==================================================================
\section*{Acknowledgements}
%%% ==================================================================

The First-G Law was originally stated (in a weaker semiprime-only form)
in WP34 \cite{WP34} of the TIG / CK research program maintained by 7Site
LLC and the Coherence Keeper community at
\href{https://github.com/TiredofSleep/ck}{github.com/TiredofSleep/ck}. The
present paper upgrades that statement to the Coprimality-Partition
Localization form above, which is the version needed as a foundational
lemma in the companion papers \cite{WP35,WP101} on composition-table
dynamics and non-associativity rates. We thank the Coherence Keeper
collaborators (Ben Mayes, H.~J.~Johnson) for discussions on the sieve
reorganization.

%%% ==================================================================
\begin{thebibliography}{99}
%%% ==================================================================

\bibitem{Apostol1976}
T.~M. Apostol, \emph{Introduction to Analytic Number Theory}, Undergraduate
Texts in Mathematics, Springer-Verlag, New York, 1976.

\bibitem{HardyWright2008}
G.~H. Hardy and E.~M. Wright, \emph{An Introduction to the Theory of
Numbers}, 6th ed., revised by D.~R. Heath-Brown and J.~H. Silverman, Oxford
University Press, Oxford, 2008.

\bibitem{IrelandRosen1990}
K.~F. Ireland and M.~I. Rosen, \emph{A Classical Introduction to Modern
Number Theory}, 2nd ed., Graduate Texts in Mathematics 84, Springer-Verlag,
New York, 1990.

\bibitem{Lang2002}
S.~Lang, \emph{Algebra}, 3rd ed., Graduate Texts in Mathematics 211,
Springer-Verlag, New York, 2002. (Chinese Remainder Theorem and
idempotents in semisimple rings: Chapter II, \S 2.)

\bibitem{Montgomery1973}
H.~L. Montgomery, \emph{The pair correlation of zeros of the zeta function},
in Analytic Number Theory (St.~Louis, MO, 1972), Proc.~Sympos.~Pure Math.
\textbf{24}, American Mathematical Society, 1973, pp.~181--193.

\bibitem{Shannon1949}
C.~E. Shannon, \emph{Communication in the presence of noise}, Proc. IRE
\textbf{37}(1) (1949), 10--21.

\bibitem{WP34}
B.~R. Sanders, \emph{The First-G Law and Prime-Forced Dispersion}, working
paper WP34 (2026), 7Site LLC research register;
DOI \texttt{10.5281/zenodo.18852047}.

\bibitem{WP35}
B.~R. Sanders, C.~A. Luther, M.~Gish, \emph{The Prime Phase Transition:
Harmonic Pre-Echo, Zero-Width Gates, and the Geometry of RSA Security},
working paper WP35 (2026), 7Site LLC research register;
DOI \texttt{10.5281/zenodo.18852047}.

\bibitem{WP101}
B.~R. Sanders, M.~Gish, C.~A. Luther, H.~J. Johnson, \emph{Non-Associativity
Decay in Binary Composition Tables over $\Z/N\Z$}, working paper WP101
(2026), 7Site LLC research register; DOI \texttt{10.5281/zenodo.18852047}.

\bibitem{proof}
B.~R. Sanders and collaborators, verification script
\texttt{proof\_first\_g\_event.py}, Sprint 35 bundle
(2026-04-19), Coherence Keeper repository,
\href{https://github.com/TiredofSleep/ck}{github.com/TiredofSleep/ck}.

\end{thebibliography}

\end{document}
